\session{7}{10/24}{2限}{2}{ケーススタディ：企業のAI活用}

\goals{
  \item 企業のAI導入事例を、目的・対象業務・データ・効果の測り方・つまずきの要因の観点で読み解ける。
  \item PoC（概念実証）の意味と、PoC止まりの問題を説明できる。
  \item 複数の事例から共通する成功要因と失敗要因を抽出し、自社への応用を考えられる。
}

\guidesec{解説}

\topic{事例を読み解く枠組み}
企業のAI活用事例は、メディアやベンダーの資料に数多く公開されています。しかし、多くは成功を強調した短い記事であり、
そのまま真似しても同じ結果にはなりません。事例を読むときは、次の観点で情報を整理します。
\par\noindent\begin{gtabh}{colspec={Q[l,wd=26mm] X[l]}}
  観点 & 確認すること \\
  目的 & 解決したかった経営上の課題は何か。「AIを使うこと」が目的になっていないか。 \\
  対象業務 & どの業務の、どの作業にAIを使ったか。業務全体か一部か。 \\
  データ & どんなデータを、どれだけ、どこから集めたか。データの整備にどれだけかかったか。 \\
  技術と体制 & 自社開発か外部サービスか。誰が推進し、現場はどう関わったか。 \\
  効果の測り方 & 導入前に決めた指標（KPI）は何か。導入前後で比較しているか。効果の数値に出典があるか。 \\
  つまずき & 途中で何が問題になり、どう対処したか。記事に書かれていないことは何か。 \\
  前提条件 & 業種、規模、データの有無、組織文化など、自社と異なる点。 \\
\end{gtabh}

\topic{効果の測り方と PoC}
AI導入の効果を示すには、導入前に目的と評価指標（KPI：Key Performance Indicator）を決め、導入前後で比較できるようにしておく必要があります。
KPI の例は、処理時間、対応件数、誤り率、顧客満足度、売上、コストなどです。導入後に「なんとなく便利になった」と聞いても、効果は示せません。

本格導入の前に、小規模で技術的な実現可能性と効果を検証する取り組みを \textbf{PoC（Proof of Concept、概念実証）} と呼びます。
PoC は有効な方法ですが、多くの企業で「PoC は成功したが本番に進まない」「PoC を繰り返すだけで終わる」という
\textbf{PoC止まり}が課題になっています。原因は、PoC の段階で「どの結果が出たら本番に進むか」という判断基準と、
本番展開に必要な予算・体制・データ整備の見通しを決めていないことです。

\topic{成功要因と失敗要因}
公開事例と調査研究から、共通して指摘される要因をまとめます。
\par\noindent\begin{gtabh}{colspec={X[l] X[l]}}
  成功しやすい条件 & 失敗しやすい条件 \\
  解決したい課題と KPI が明確 & 目的が曖昧で、技術先行（「AIで何かできないか」） \\
  経営層が目的を示し、推進に責任を持つ & 現場やベンダーに丸投げ \\
  使えるデータがあり、整備に投資している & データが散在・欠損していて、整備を軽視 \\
  小さく始めて効果を測り、段階的に広げる & 最初から大規模に導入し、失敗が見えにくい \\
  現場が設計に参加し、業務を見直している & 既存の業務にAIを「貼り付ける」だけ \\
  人の確認と例外処理を設計している & AIの出力を無条件に使い、誤りで信頼を失う \\
  効果を継続的に測り、改善している & 導入が目的化し、導入後に放置 \\
\end{gtabh}

\topic{代表的な活用領域}
よく見られる活用領域を挙げます。事例を探すときの手がかりにしてください。
\begin{itemize}
  \item \textbf{顧客対応}　問い合わせ対応のチャットボット、オペレーターへの回答候補の提示、通話の要約。
  \item \textbf{文書業務}　契約書・報告書の下書きと要点抽出、社内規程の検索と回答、翻訳。
  \item \textbf{需要予測と在庫}　販売データと外部要因（天候、イベント）からの需要予測、発注の自動化。
  \item \textbf{製造・品質}　画像による不良検知、設備の異常予知。
  \item \textbf{ソフトウェア開発}　コードの生成・補完・レビュー、テストの自動化。
  \item \textbf{マーケティング}　顧客ごとの提案（パーソナライズ）、広告文の生成と検証。
\end{itemize}

\topic{キーワード}
\par\noindent\begin{keywords}{}
  KPI & 目的の達成度を測る評価指標。導入前に決め、導入前後で比較する。 \\
  PoC & 本格導入の前に小規模で実現可能性や効果を検証すること。 \\
  PoC止まり & PoC を繰り返すだけで本番展開に進まない状態。判断基準と体制の不在が原因。 \\
  技術先行 & 解決したい課題より先に技術の導入を決めてしまうこと。典型的な失敗要因。 \\
  前提条件 & 事例を自社に応用するときに確認する、業種・規模・データ・組織文化の違い。 \\
\end{keywords}

\guidesec{演習課題}

\exercise{1}{AI演習：複数の事例をAIで要約・比較し、共通する成功要因を討議する}
\instr{企業のAI活用事例を3件、公開情報（企業の公式サイト、省庁の事例集、報道）から選び、それぞれの記事の本文を次のプロンプトに貼って要約・比較させてください。比較表を見て、3件に共通する成功要因と、記事に書かれていない「つまずき」の候補を考え、グループで討議してください。}
\begin{promptbox}[事例の要約と比較]
次の3つの企業のAI活用事例について、「目的」「対象業務」「使ったデータ」「技術と体制（自社開発か外部サービスか）」「効果の測り方とKPI」「つまずきと対処」「記事に書かれていない重要な情報」の7項目で表に整理してください。記事に書かれていない項目は「記載なし」としてください。推測で埋めないでください。
（事例1の本文）（事例2の本文）（事例3の本文）
\end{promptbox}

\exercise{2}{失敗事例を分析する}
\instr{次の架空の事例を読み、失敗の要因を3つ指摘し、それぞれ「何をしていれば防げたか」を書いてください。}
\begin{point}{事例：ある地方の食品卸売業（従業員120人）}
社長が展示会で生成AIのデモを見て、「我が社も使うべきだ」と決断。情報システム担当の1名に導入を指示し、
3か月後に全社員にチャット型AIの利用を開放した。説明会は30分のオンライン1回のみ。半年後、
利用しているのは数名で、営業担当の一人が顧客の取引条件を個人のAIアカウントに入力していたことが発覚。
社長は「効果がなかった」として契約を打ち切った。
\end{point}

\exercise{3}{自社への応用を考える}
\instr{演習1で選んだ事例のうち1件について、自分が働きたい業界または身近な企業に応用する場合、前提条件の違い（業種、規模、データ、組織文化）を表にし、「そのまま使える点」「変える必要がある点」「応用できない点」に分けてください。}

\guidesec{模範解答}

\begin{answer}{演習1}
比較表の例（3件の記事を要約したもの）：
\par\noindent\begin{gtabh}{colspec={Q[l,wd=22mm] X[l] X[l] X[l]}}
  項目 & コールセンター（保険） & 社内文書検索（製造） & 需要予測（小売） \\
  目的 & 応対時間の短縮と新人の早期戦力化 & 規程・技術文書を探す時間の削減 & 廃棄ロスの削減 \\
  対象業務 & 通話中の回答候補の提示、応対後の要約 & 社内文書への質問応答 & 日配品の発注量の決定 \\
  データ & 過去の応対記録、FAQ & 社内規程、設計文書 & POS データ、天候、イベント \\
  体制 & 外部サービス＋自社で FAQ 整備 & 外部の LLM＋自社開発の検索 & ベンダーと共同開発 \\
  KPI & 平均応対時間、新人の独り立ち期間 & 記載なし（「好評」とのみ） & 廃棄率、欠品率 \\
  つまずき & 回答候補の誤りで現場が使わなくなった時期があり、FAQ を整備し直した & 記載なし & 記載なし \\
  書かれていないこと & 導入費用、効果測定の期間 & 効果の数値、利用率 & 予測が外れた場合の運用 \\
\end{gtabh}
共通する成功要因（討議の結論の例）：目的が業務の数値（時間、廃棄率）で定義されている。自社のデータ（FAQ、文書、POS）を
整備する作業に投資している。外部サービスを使いつつ、自社固有の部分（FAQ、検索、運用）は自社で設計している。
書かれていない「つまずき」の候補：効果の数値の測定期間、現場の利用率、誤りが起きたときの対応。これらは事例の企業に直接聞くか、
講演資料や学会発表を探すと見つかることがあります。
\end{answer}

\begin{answer}{演習2}
\begin{enumerate}[label=\arabic*.,leftmargin=2em]
  \item 目的が曖昧で技術先行：「使うべきだ」から始まり、解決したい課題と KPI がない。防ぐには、導入前に「受発注の入力時間を減らす」など課題を特定し、指標を決める。
  \item 体制と教育の不足：担当1名への丸投げと、30分の説明会のみ。防ぐには、経営層が推進責任を持ち、部門ごとの推進担当と実践的な研修、相談窓口を用意する。
  \item ルールと公式ツールの不在：使ってよい情報の範囲が示されず、個人アカウントで機密情報が入力された（シャドーAI）。防ぐには、法人契約のツールを提供し、入力してよい情報のルールを導入と同時に周知する。
\end{enumerate}
加えて、「効果がなかった」という判断自体が、効果を測る指標がなかったために根拠を欠いています。
\end{answer}

\begin{answer}{演習3}
例：需要予測（大手小売）を、地方の食品スーパー（3店舗）に応用する
\par\noindent\begin{gtabh}{colspec={Q[l,wd=20mm] X[l] X[l]}}
  前提条件 & 事例の企業 & 応用先 \\
  規模 & 数百店舗、専任のデータ部門 & 3店舗、専任者なし \\
  データ & 全店の POS が統合されている & POS はあるが店舗ごとに別。天候データは未取得 \\
  組織文化 & 本部主導の発注 & 各店の担当者の経験による発注 \\
\end{gtabh}
\begin{itemize}
  \item そのまま使える点：廃棄率と欠品率を KPI にする考え方。日配品から始める対象の絞り込み。
  \item 変える必要がある点：自社開発ではなく、POS ベンダーが提供する予測機能（SaaS）を使う。担当者の経験を否定せず、予測を「参考値」として示すことから始める。
  \item 応用できない点：店舗横断の在庫融通や、大規模なデータ基盤の構築。
\end{itemize}
\end{answer}

\guidesec{出席確認クイズの解説}
講義サイトの出席確認ページ（第7回）の5問です。
% 出席確認クイズ 第07回　ケーススタディ：企業のAI活用
%   attendance/attendance07.html から生成（解説は本ガイド用に加筆）。

\quizq{1}{AI導入の効果を測る際に、まず必要なことはどれでしょうか?}%
  {効果測定は行わない}%
  {導入費用が高いほど効果が高いとみなす}%
  {\correct{導入前に目的とKPI(評価指標)を決めておく}}%
  {導入後になんとなく満足度を聞く}%
  {正解は (c)。効果を測るには、導入前に目的と評価指標（KPI）を定め、導入前後で比較できるようにしておく必要があります。(d) のように後からなんとなく満足度を聞いても比較できません。(b) 費用と効果は別のものです。}

\quizq{2}{「PoC(概念実証)」の説明として適切なものはどれでしょうか?}%
  {契約書の一種}%
  {本番システムの名称}%
  {\correct{本格導入の前に小規模で実現可能性や効果を検証すること}}%
  {AIの学習アルゴリズムの名前}%
  {正解は (c)。PoC（Proof of Concept、概念実証）は、本格導入の前に小規模で技術的な実現可能性や効果を確かめる取り組みです。PoC が次に進まず繰り返される「PoC止まり」が多くの企業で課題になっており、PoC の段階で本番展開の条件を決めておくことが重要です。}

\quizq{3}{AI導入が失敗しやすい典型的な要因として適切なものはどれでしょうか?}%
  {データの整備をしすぎる}%
  {現場を巻き込みすぎる}%
  {\correct{目的が曖昧なまま技術先行で導入する}}%
  {目的が明確すぎる}%
  {正解は (c)。「AIを使うこと」自体が目的化し、解決したい課題が曖昧なまま技術先行で導入するのが典型的な失敗要因です。(a)(b) は「やりすぎ」を装った選択肢ですが、データの整備と現場の巻き込みはむしろ成功要因です。}

\quizq{4}{顧客対応でのAI活用例として一般的なものはどれでしょうか?}%
  {商品の物理的な配送}%
  {電力の供給}%
  {店舗の清掃}%
  {\correct{チャットボットによる問い合わせ対応}}%
  {正解は (d)。問い合わせ対応のチャットボットは、顧客対応における AI 活用の代表的な例です。(a)(b)(c) は物理的な作業で、生成AIが直接担う領域ではありません。チャットボットでも、回答の正確さと人への引き継ぎの設計が成否を分けます。}

\quizq{5}{他社の事例を読み解く際の姿勢として適切なものはどれでしょうか?}%
  {有名企業の事例だけを見ればよい}%
  {成功事例をそのまま真似すれば必ず成功する}%
  {失敗事例には見る価値がない}%
  {\correct{業種・規模・データの条件が自社と異なる点を踏まえて比較する}}%
  {正解は (d)。事例は業種・規模・データの条件などの前提が異なるため、自社との違いを踏まえて何が応用できるかを考えることが大切です。(b) 成功事例の模倣が必ず成功するわけではなく、(c) 失敗事例からはつまずきの要因を学べます。(a) 規模の近い企業の事例の方が参考になることも多いです。}

